% 379_Lien_Plenum_Photon_En_ch.tex
% Johann Pascher, 29 Sept. 2026
% The acoustic plenum and the photon: Lien (2026) in the light of
% Docs. 267, 290, 340, A265 and R128

\providecommand{\BM}{\textbf{[B]}}
\providecommand{\KM}{\textbf{[K]}}
\providecommand{\SM}{\textbf{[S]}}
\providecommand{\XM}{\textbf{[X]}}
\providecommand{\QM}{\textbf{[Q]}}

\chapter*{The Acoustic Plenum and the Photon}
\addcontentsline{toc}{chapter}{The Acoustic Plenum and the Photon}

\begin{abstract}
C.~C.~Lien (Zenodo, DOI 10.5281/zenodo.21800087, 8~September 2026) interprets
space as a compressible superfluid, derives the electron mass from the
temperature of the cosmic microwave background, adopts a mass formula for the
muon and tau, equates the neutrino mass with $k_BT_\text{CMB}$ and describes the
photon as a longitudinal shockwave with a transverse wake. This document
examines the paper as a case study with a verification script (20/20). The
numerical values are reproducible \KM. None of the load-bearing claims
survives the check: the electron hit lies 88 standard deviations from the
measured value and is not significant in the null test, the lepton formula is
known, the neutrino mass is incompatible with oscillation data, the redshift
would be frequency-dependent, and the photon would have three instead of two
polarisations \XM. The cross-check shows how FFGFT proceeds at the same
points, without repeating those contents.
\end{abstract}

% ============================================================
\section*{Status markers}
% ============================================================
\begin{center}
\begin{tabular}{@{}lp{10cm}@{}}
\textbf{[B]} & Algebraically proved. \\[3pt]
\textbf{[K]} & Numerically confirmed (verification script). \\[3pt]
\textbf{[S]} & Structurally plausible, open. \\[3pt]
\textbf{[X]} & Examined and rejected. \\[3pt]
\textbf{[Q]} & Quoted verbatim from the source (Lien 2026). \\
\end{tabular}
\end{center}

% ============================================================
\section{The source}
% ============================================================

Lien replaces empty space with a medium:
\begin{quote}\QM
\enquote{By treating the spatial manifold as a continuous compressible
superfluid, we derive the foundational particle phase-states (electron, muon,
tau, and neutrino) from the Cosmic Microwave Background (CMB) acoustic
baseline.} (Conclusion)
\end{quote}
The background is stationary, and redshift is not expansion:
\begin{quote}\QM
\enquote{Cosmological redshift is therefore not temporal expansion; it is
Acoustic Pulse Broadening -- the observation that waves of higher energetic
frequencies decay faster than those of a lower energetic frequency when
traversing the same medium.} (Section II.A)
\end{quote}
The photon is not a particle:
\begin{quote}\QM
\enquote{Photonic emission is therefore not a discrete point-particle
phenomenon. It is a primary longitudinal acoustic shockwave [\ldots] which
structurally interacts with the azimuthal shear of its source to continuously
generate an orthogonal transverse spiral wake.} (Section V)
\end{quote}
The central formulas are
\begin{align}
f_e &= 2\pi\,(\alpha^{-1})^4\,\frac{k_BT_\text{CMB}}{h}, \label{eq:lien-fe}\\
m_n &= m_e\Bigl(1+\tfrac{3}{2}\,\alpha^{-1}\sum_{k=0}^{n}k^4\Bigr), \label{eq:lien-mn}\\
m_\nu &= k_BT_\text{CMB}/c^2. \label{eq:lien-mnu}
\end{align}

% ============================================================
\section{Recalculation}
% ============================================================

All six numerical values are reproducible \KM:
$f_\text{CMB}=5.679\times10^{10}$\,Hz, $f_e=1.2583\times10^{20}$\,Hz
against the Compton frequency $1.2356\times10^{20}$\,Hz ($+1.84\,\%$),
$m_\mu/m_e=206.55$ ($-0.10\,\%$), $m_\tau/m_e=3495.4$ ($+0.52\,\%$) and
$m_\nu=0.235$\,meV. There are no arithmetic errors. The only question is
what the numbers mean.

% ============================================================
\section{Significance}
% ============================================================

\paragraph{The electron hit is a miss.}
The ratio $m_ec^2/(k_BT_\text{CMB})=2.176\times10^9$ lies $1.84\,\%$ below
$2\pi\alpha^{-4}=2.216\times10^9$. $T_\text{CMB}$ is known to $0.02\,\%$; the
deviation therefore amounts to about 88 standard deviations \KM. Lien
interprets it as \enquote{dynamic acoustic drag} but does not calculate this
drag. A friction named after the fact is not a correction but a second free
parameter.

\paragraph{Null test.}
The formula family $a\,\pi^b\alpha^{-k}$ with thirteen simple fractions $a$,
$b\in\{-3,\dots,3\}$ and $k\in\{1,\dots,6\}$ comprises 546 formulas. Against
random target values between $10^9$ and $10^{10}$, a hit within $1.84\,\%$
succeeds in 61\,\% of cases \KM. The hit is therefore not significant. FFGFT
applies the same standard to its own grid findings (Doc.~342, 343 §G$'$):
what chance delivers equally well is not counted.

\paragraph{The lepton formula is known.}
Eq.~\eqref{eq:lien-mn} is the formula of A.~O.~Barut (Phys.\ Rev.\ Lett.\
42, 1251, 1979) without change \BM. The accuracy for muon and tau is Barut's
result, not that of the plenum. For $n=3$ the same formula would give a
charged lepton at $10.3$\,GeV \KM, which is experimentally excluded; Lien's
justification via Chirikov overlap remains qualitative. For the number of
generations in FFGFT see Doc.~317 and 316.

% ============================================================
\section{Redshift}
% ============================================================

\paragraph{Damping changes the amplitude, not the frequency.}
Stokes damping $\alpha\propto f^2$ reduces the amplitude of a wave; its
frequency is preserved. Numerically, a damped 50\,Hz oscillation keeps its
spectral maximum at 50\,Hz \KM. No redshift arises in this way.

\paragraph{A frequency-dependent redshift is excluded.}
If $z$ were tied to damping $\propto f^2$, the 21\,cm line and Lyman-$\alpha$
would differ by a factor of $3\times10^{12}$ \KM. As measured, $z$ agrees
between the two to about $10^{-5}$. FFGFT went through this error itself:
older corpus documents contained a chromatic redshift, which A265 marks as
superseded.

\paragraph{Light-curve stretching.}
Lien mentions the $(1+z)$ stretching of supernova light curves only as an
argument of standard cosmology, without explaining it. A static picture must
contain it. In FFGFT it follows from $\tilde T\cdot m=1$: wavelength and clock
rate are booked together, and the stretching is exactly $(1+z)$ for all $z$
(Doc.~312, A265, R128) \BM. By Doc.~267 all frequency and clock-rate
signatures are degenerate between a static picture and expansion; a static
background is therefore admissible, but only if it is achromatic. Lien's
plenum is not \XM.

% ============================================================
\section{The photon}
% ============================================================

In vacuum, light has exactly two transverse polarisations and no
longitudinal component. A massless $U(1)$ gauge field has four components,
two of which are removed by gauge freedom and the constraint \BM. Lien's
picture of a longitudinal primary wave plus a transverse wake would have
three degrees of freedom \XM. The baroclinic torque $\nabla\rho\times\nabla P$
generates a vorticity component in the $z$ direction, but no wave with two
polarisation states. Photon counting (photoelectric effect, antibunching) is
not addressed in the paper.

The argument that a particle picture would have to break up into a
\enquote{quantised spray} in the far field confuses the wavefront of a
classical field with the counting statistics of its quanta. In FFGFT the
photon is exactly massless, the pure phase and connection anchor, and
$E_\text{bit}=\hbar c/L$ is a photon energy (Doc.~290). No medium is needed
for this.

% ============================================================
\section{Neutrino}
% ============================================================

A single mass of $0.235$\,meV cannot produce the atmospheric mass splitting
$\Delta m^2_\text{atm}\approx2.5\times10^{-3}\,\text{eV}^2$. For that, at least
one state must weigh about 50\,meV \KM. Equating the neutrino mass with
$k_BT_\text{CMB}$ is therefore excluded \XM. The FFGFT hierarchy of Doc.~340
satisfies this condition with its heaviest state at $49.96$\,meV \KM.

% ============================================================
\section{The shared anchor}
% ============================================================

It is striking that Lien touches the same number that is the SI anchor of
the cosmic sector in FFGFT: $m_ec^2/(k_BT_\text{CMB})=2.176\times10^9$ (A265).
FFGFT carries this ratio explicitly as an anchor, not as a derivation; the
exponent $41/4$ of the bridge formula is open (P20), and the cosmic sector
is deliberately bracketed (P39). Lien's candidate $2\pi\alpha^{-4}$ misses the
number by $1.84\,\%$ and cannot be distinguished from chance in the null
test. It therefore does not close the gap; it shows how easily
approximations can be found at this point \SM.

% ============================================================
\section{What agrees and what does not}
% ============================================================

The paper shares two starting points with FFGFT: a stationary instead of an
expanding background, and the rejection of point-like singularities. In
FFGFT the smallest length $L_0=\xi\,\ell_P$ secures finiteness (Doc.~180,
R90); in Lien's work a density maximum $\rho_E$ does this. That maximum,
however, is circular, because the classical electron radius
$r_e=\alpha\hbar/(m_ec)$ already contains the electron mass \KM.

The difference lies in the method. Lien introduces a medium with material
constants (damping, viscosity, drag) and explains deviations after the fact
through these constants. FFGFT posits a compact geometry with one parameter
$\xi$, derives numbers and checks them against fixed criteria. Where a
deviation remains, it is carried as an open bridge, not explained away.

% ============================================================
\section{Summary}
% ============================================================

\begin{center}
\begin{tabular}{>{\raggedright\arraybackslash}p{4.2cm}>{\raggedright\arraybackslash}p{4.8cm}>{\raggedright\arraybackslash}p{2.2cm}l}
\toprule
Claim (Lien) & Finding & FFGFT & Status \\
\midrule
$f_e=2\pi\alpha^{-4}f_\text{CMB}$ & $+1.84\,\%$, 88\,$\sigma$, null test 61\,\% & A265, P20 & \XM \\
Lepton formula Eq.~\eqref{eq:lien-mn} & Barut 1979, $n=3$ excluded & 316, 317 & \XM \\
$m_\nu=k_BT_\text{CMB}$ & incompatible with $\Delta m^2_\text{atm}$ & 340 & \XM \\
Redshift by damping & frequency-dependent, leaves $f$ unchanged & 267, 312, R128 & \XM \\
Photon longitudinal + transverse & three instead of two degrees of freedom & 290 & \XM \\
Stationary background & admissible if achromatic & 267, R128 & \SM \\
\bottomrule
\end{tabular}
\end{center}

The paper combines known formulas, a friction named after the fact and an
intuitive medium into a picture that appears closed. The recalculation shows
that the numbers are correct, but that the interpretation fails against
measured data at every load-bearing point. The usable core, a static
background without singularities, is already realised differently and
achromatically in FFGFT.

\paragraph{Verification script.} \texttt{2/python/Dok379\_Skripte/pruef\_379\_lien\_photon.py}
(20/20): recalculation (A), significance and null test (B), measured data
(C), FFGFT cross-check (D). Two inputs are literature values and marked as
such: the agreement of $z$ between 21\,cm and Lyman-$\alpha$, and the DES value
$b=1.003\pm0.011$.

\paragraph{Source document.} Lien, C.~C.: \textit{Electromagnetic Illusion:
Photonic Emission as Orthogonal Acoustic Shockwaves in a Compressible
Superfluid Plenum.} Zenodo (2026), DOI 10.5281/zenodo.21800087.
Barut, A.~O.: \textit{Lepton Mass Formula.} Phys.\ Rev.\ Lett.\
\textbf{42}, 1251 (1979).
